% =============================================================================
%  Supplementary Section S6: dosing interval specification.
%  Order: the trough of a dose train; the trough condition and the crossing half-life;
%  the ceiling; crossings over the carnitine constants; evidence on the half-life;
%  uncertainty in the half-life; a threshold that allows for illness; estimating theta
%  instead; reading the half-life from the dated record.
%  Sources: audit/brief_2026-10-01/pi_note/pi_note.tex (Result 4), figure_guide.tex
%  (Figure 4), out/t44, t45, t46, t48, t50, t52, t53, t54; REPORT.md R-J to R-S.
% =============================================================================
\section{Dosing interval specification}
\label{si:s6}

This section derives the condition behind the dosing specification of the main text and
reports the numbers on which it rests. All results use the carnitine model at its baseline
load of 600\um\ (Section~\ref{si:s1-carn}). The threshold is the intrinsic critical capacity
$\thci$, the least enzyme capacity at which a patient at that load still has a stable state
(Eq.~\eqref{seq:s1-thcint}); at the nominal carnitine constants it is 150.8\um, about 5\% of
normal. Residual capacity $\theta$ (normal 3000\um) enters only as an additive term, and the
question throughout is when it can be left out of the decision.
% src: audit/brief_2026-10-01/out/t53_crossing_shape.txt:5 (theta_crit_int 150.775); manuscript/main.tex:386-389

% -----------------------------------------------------------------------------
\subsection{The trough of a dose train}
\label{si:s6:trough}

Let infusions be given every $\tau$ days, at $t_k=k\tau$. We hold the amount delivered per day
equal to that of the fortnightly schedule, so that the jump in delivered capacity at each
infusion is $D_\tau=D\,\tau/14$, with $D=6000$\um\ the assumed fortnightly jump; a 21-day
infusion is then 1.5 times a 14-day one. By Eq.~\eqref{seq:s1-delta}, just before the
infusion at $t=n\tau$, after $n$ infusions at $t_0,\dots,t_{n-1}$, the delivered increment is
\begin{equation}
\label{seq:s6-partial}
\Delta\bigl((n\tau)^{-}\bigr)\;=\;D_\tau\sum_{j=1}^{n}x^{j}\;=\;D_\tau\,\frac{x\,(1-x^{n})}{1-x},
\qquad x=e^{-\lambda\tau},\qquad \lambda=\frac{\ln2}{t_{1/2}} .
\end{equation}
The factor $x^{n}$ shrinks by $x$ with every infusion ($x=1/16$ at 14 days and a 3.5-day
half-life), so the train settles within a few infusions to a periodic pattern with trough and
peak
\begin{equation}
\label{seq:s6-trough}
\dtr(\tau)\;=\;D\,\frac{\tau}{14}\,\frac{e^{-\lambda\tau}}{1-e^{-\lambda\tau}},\qquad
\Delta_{\mathrm{pk}}(\tau)\;=\;D\,\frac{\tau}{14}\,\frac{1}{1-e^{-\lambda\tau}}
\;=\;\dtr(\tau)\,e^{\lambda\tau}.
\end{equation}
The trough is reached just before each infusion and the peak just after it. Their ratio,
$e^{\lambda\tau}=2^{\tau/t_{1/2}}$, depends only on how many half-lives the interval spans. The
mean over a cycle, $D_\tau/(\lambda\tau)=D/(14\lambda)$, is the same at every interval, so
lengthening the interval at the same daily amount lowers the trough and raises the peak without
changing the average delivered enzyme. At a 3.5-day half-life the fortnightly train has
$\dtr(14)=400$\um\ and $\Delta_{\mathrm{pk}}(14)=6400$\um, which reproduce to the decimal the
dose train integrated numerically in our code.
% src: audit/brief_2026-10-01/out/t44_interval_decision.txt:4-5 (400.0, 6400.0 reproduce the repository's train); manuscript/main.tex:741-748 (eq:trough); REPORT.md:875 (reproduce delta() numerically to the decimal)

The same-amount-per-day convention is a choice. If every infusion kept the size $D$ whatever
the interval, the trough at 15 days would be 324\um\ instead of 347\um\ at a 3.5-day half-life,
and lower still at longer intervals; at shorter intervals the fixed size gives the larger trough
(2000 against 1000\um\ at 7 days).
% src: audit/brief_2026-10-01/out/t44_interval_decision.txt:17-21 (rate vs fixed: 7 d 1000.0/2000.0; 15 d 347.4/324.3)

Because $\dtr$ falls steeply as $\lambda$ rises, the trough is very sensitive to the half-life
(Supplementary Table~\ref{stab:s6-trough}). Over half-lives from 2 to 7 days the fortnightly
trough runs from 47 to 2000\um, a factor of 42, and at a 3.5-day half-life the 21-day trough
(143\um) already lies below $\thci$.
% src: audit/brief_2026-10-01/out/t44_interval_decision.txt:140-147 (47.2 -> 2000.0, factor 42); out/t48_lambda_literature.txt:40 (Dtr21 143 at 3.5 d)

\begin{table}[tbp]
\centering
\small
\caption{\label{stab:s6-trough}Delivered capacity left at the trough, $\dtr(\tau)$ (\um), at
the assumed dose ($D=6000$\um\ per fortnight, scaled to the same amount per day at each
interval), from Eq.~\eqref{seq:s6-trough}. The threshold it must clear for a patient with no
enzyme of their own is $\thci=150.8$\um\ (nominal carnitine constants). Rows are the ends of the
plausible range of the hepatic half-life (3 and 6\,d), the value used in the identifiability and
sampling analyses (3.5\,d), the human fibroblast value (3.6\,d), the likely hepatic value
(4.3\,d) and an intermediate value (5\,d).}
\begin{tabular}{@{}llccc@{}}
\toprule
$t_{1/2}$ (d) & role & 14 days & 15 days & 21 days \\
\midrule
3.0 & lower end of the plausible range & 246 & 207 & 71 \\
3.5 & identifiability and sampling analyses & 400 & 347 & 143 \\
3.6 & human fibroblasts~\cite{chandler1985} & 434 & 379 & 161 \\
4.3 & likely hepatic value & 702 & 629 & 316 \\
5.0 & & 1006 & 918 & 518 \\
6.0 & upper end of the plausible range & 1485 & 1380 & 873 \\
\bottomrule
\end{tabular}
\end{table}
% src: audit/brief_2026-10-01/out/t44_interval_decision.txt:142-144 (Dtr14 245.9, 400.0, 1006.0); out/t48_lambda_literature.txt:40-41 (3.5 d 400/347/143; 3.6 d 434/379/161); out/t50_hepatic_mito.txt:19-23 (Dtr21 71, 161, 316, 518, 873); pi_note/pi_note.tex:431 (Dtr14 702 at 4.3 d). Cells 3.0 d/15 d (207), 4.3 d/15 d (629), 5.0 d/15 d (918), 6.0 d/14 d (1485) and 6.0 d/15 d (1380) are evaluations of eq. seq:s6-trough at the stated inputs (not printed by a script); every printed cell is reproduced by the same evaluation.

% -----------------------------------------------------------------------------
\subsection{The trough condition and the crossing half-life}
\label{si:s6:crossing}

At constant load a stable state exists while the enzyme capacity exceeds $\thci$
(Section~\ref{si:s1-carn}). Under therapy the capacity is $\theta+\Delta(t)$, which is lowest at
the trough, so we require
\begin{equation}
\label{seq:s6-condition}
\theta+\dtr(\tau)\;>\;\thci .
\end{equation}
A patient who meets Eq.~\eqref{seq:s6-condition} has a stable state at baseline load at every
point of the dose cycle. The condition is quasi-static: it asks whether a stable state would
exist if the capacity reached at the trough were held, and it does not follow the metabolite
through the cycle. Supplementary Section~\ref{si:s7} integrates the two-state system through
late infusions and compares the two criteria.

Because $\theta\ge0$, an interval keeps every patient maintainable, whatever $\theta$, when the
delivered enzyme alone clears the threshold, $\dtr(\tau)>\thci$. The margin of a patient with no
enzyme of their own, $\mu(\tau)=\dtr(\tau)-\thci$, is the quantity plotted in the main text's
margin figure. When it is negative, the interval is safe only for patients whose own capacity
exceeds $\theta_{\mathrm{need}}(\tau)=\thci-\dtr(\tau)$, and the decision then needs $\theta$,
which the routine record does not supply in severe PA (main text and Supplementary
Section~\ref{si:s4}).

For fixed $\tau$ and $D$, $\dtr$ increases strictly with the half-life, so there is exactly one
half-life, the crossing $t_{1/2}^{\ast}(\tau)$, at which $\dtr(\tau)=\thci$. Writing
$x=e^{-\lambda\tau}$, Eq.~\eqref{seq:s6-trough} gives $x=\thci/(\thci+D\tau/14)$, and hence
\begin{equation}
\label{seq:s6-cross}
t_{1/2}^{\ast}(\tau)\;=\;\frac{\tau\,\ln2}{\ln\!\bigl(1+D\tau/(14\,\thci)\bigr)} .
\end{equation}
No iteration is needed. Above the crossing the interval is safe without knowing $\theta$; below
it, it is not. The crossing depends on the dose and the threshold only through their ratio
$D/\thci$, and it lengthens as the dose falls or the threshold rises. At $D=6000$\um\ and the
nominal constants it is 2.62, 2.75 and 3.55 days for intervals of 14, 15 and 21 days.
% src: audit/brief_2026-10-01/out/t53_crossing_shape.txt:5-6 (2.617, 2.754, 3.5452 d); eq. seq:s6-cross is the closed form of the same root and reproduces these values

At the half-life of 3.5 days used in the identifiability and sampling analyses, the 15-day
trough (347\um) clears $\thci$ by 197\um, so moving from 14 to 15 days would need no
measurement of $\theta$. The 21-day trough (143\um) falls 7.9\um\ short, so the 21-day interval
would be safe only for patients with at least 7.9\um\ of their own (0.26\% of normal), a level
that the record cannot confirm.
% src: audit/brief_2026-10-01/out/t46_timing.txt:5-6 (mu_long = theta_base - 7.9), :64-65 (347.4 > 150.8, margin +196.6); out/t54_dose_timing.txt:94 (theta_need 7.9, 0.26 %)

% -----------------------------------------------------------------------------
\subsection{The ceiling, and why the assumed dose breaks it}
\label{si:s6:ceiling}

The upper edge of the therapeutic window is the ceiling of Section~\ref{si:s1-eq}: just after an
infusion, the patient's own capacity plus the delivered enzyme should not exceed normal
capacity. We use it as a plausibility condition of the model, not as a measured toxicity limit,
and set it for a patient at 5\% of normal (150\um), which caps the delivered peak at
$3000-150=2850$\um.
% src: manuscript/main.tex:117-121,753-755 (ceiling as plausibility condition; delivered peak at most 2850)

The assumed dose breaks this ceiling two- to threefold. On the fortnightly schedule at a 3.5-day
half-life it takes a patient at 5\% of normal to $(150+6400)/3000=2.18$ times normal activity
just after an infusion. At 4.3 days the delivered peak alone is about 6700\um\ (2.2 times
normal) at 14 days and about 9300\um\ (3.1 times normal) at 21 days (Supplementary
Table~\ref{stab:s6-rules}). The dose was chosen so that the marker falls by about two thirds at
peak effect (62--66\%, depending on the baseline and the peak from which the fall is read),
whereas the trial reported changes of $-29$, $-3$, $-20$, $+5$ and $-20\%$ across its five
cohorts after the third dose, with no dose trend~\cite{trial_phase12}. The ceiling and the
reported falls are two independent reasons to think the assumed dose too high.
% src: paper/sections/11_scope.tex:155-158 (2.18 times normal at 5 %; D = 2672); src/common/params.py:307-321 (D row: 62-66 %, trial falls -29, -3, -20, +5, -20 %); audit/brief_2026-10-01/pi_note/figure_guide.tex:478-484 (tops 6700 and 9300; 2.2x and 3.1x normal)

Capping the dose at the ceiling means lowering the jump at each interval and half-life until the
periodic peak equals 2850\um, $D_\tau^{\mathrm{cap}}=2850\,(1-e^{-\lambda\tau})$. On the
fortnightly schedule at 3.5 days this is $D=2672$\um, the capped dose used in Supplementary
Section~\ref{si:s7}; at 4.3 days it is about 2750\um\ for a 21-day interval, against 9000\um\ at
the assumed dose. With the peak fixed, the trough is $2850\,e^{-\lambda\tau}$ and the crossing
becomes
\begin{equation}
\label{seq:s6-capcross}
t_{1/2}^{\ast,\mathrm{cap}}(\tau)\;=\;\frac{\tau\,\ln2}{\ln\bigl(2850/\thci\bigr)} ,
\end{equation}
which no longer depends on $D$. Equivalently, at the ceiling an interval keeps every patient
maintainable only if it spans fewer than $\log_2(2850/\thci)$ half-lives, 4.24 at the nominal
constants, so that the crossings are $14/4.24=3.30$, $15/4.24=3.54$ and $21/4.24=4.95$ days.
% src: audit/brief_2026-10-01/out/t53_crossing_shape.txt:86-91 (trough = 2850 exp(-lambda tau); crossing = tau ln2/ln(2850/theta_crit); 3.30/3.54/4.95); paper/sections/11_scope.tex:158 (D = 2672); pi_note/figure_guide.tex:491-493 (about 2750 for a 21-day interval, against 9000)

At 14 days the two edges of the window can hold together: at 4.3 days a peak at the ceiling
still leaves 298\um\ at the trough. At 21 days they cannot. The delivered enzyme falls about
thirtyfold over the interval, so a peak at the ceiling leaves 97\um, and a trough of 151\um\
would need a peak of about 4450\um. The assumed dose keeps the 21-day trough high only because
its peak breaks the ceiling.
% src: audit/brief_2026-10-01/pi_note/figure_guide.tex:470-495 (298, 97, 4450; thirtyfold over 21 days)

\begin{table}[tbp]
\centering
\small
\caption{\label{stab:s6-rules}The two edges of the therapeutic window at a half-life of 4.3\,d,
in delivered capacity (\um). Peak: just after an infusion, in the periodic state. Trough: just
before the next. The lower edge asks whether the trough clears $\thci=150.8$\um\ (a patient with
no enzyme of their own; nominal carnitine constants); the upper edge asks whether the peak stays
within the ceiling of 2850\um\ set for a patient at 5\% of normal. Assumed dose: 6000\um\ per
fortnight, 9000\um\ per 21-day infusion. Capped: the largest dose whose peak stays at the
ceiling.}
\begin{tabular}{@{}llcccc@{}}
\toprule
interval & dose & peak & trough & trough above $\thci$? & peak within the ceiling? \\
\midrule
14 days & assumed & about 6700 & 702 & yes & no (2.2 times normal) \\
14 days & capped & 2850 & 298 & yes & yes \\
21 days & assumed & about 9300 & 316 & yes & no (3.1 times normal) \\
21 days & capped & 2850 & 97 & no & yes \\
\bottomrule
\end{tabular}
\end{table}
% src: audit/brief_2026-10-01/pi_note/figure_guide.tex:474-485 (table: tops 6700, 2850, 9300, 2850; bottoms 702, 298, 316, 97)

% -----------------------------------------------------------------------------
\subsection{Crossings over the unmeasured carnitine constants}
\label{si:s6:corners}

The threshold $\thci$ depends on the three shape constants of the carnitine route, $K_c$, $M_i$
and $\rho_0$, none of which has been measured (Section~\ref{si:s1-carn}). We evaluated it at the
eight corners of their assumed ranges ($K_c$ 25--800\uM, $M_i$ 150--1200\uM, $\rho_0$ 0.35--1),
re-deriving $R_0$ at each corner so that the non-enzymatic route alone still peaks at 500\um\
against the baseline load of 600\um\ (Eq.~\eqref{seq:s1-R0}). Over the corners $\thci$ spans
109.9--232.1\um, a factor of 2.1 (Supplementary Table~\ref{stab:s6-corners}). It is highest at
$K_c=25$\uM\ and $M_i=150$\uM, where the route peaks at the lowest metabolite level, and lowest
at the opposite corner; $\rho_0$ matters little (a factor of 1.02 when swept alone). We call the
corner with the highest threshold ($K_c=25$\uM, $M_i=150$\uM, $\rho_0=1$) the least favourable
for the trough condition.
% src: audit/brief_2026-10-01/out/t53_crossing_shape.txt:64-76 (corners of the assumed ranges, R0 re-derived; theta_crit 109.9..232.1, factor 2.11); results/tables/tab_C_shape.tex:5-7 (one-at-a-time factors 1.56, 1.37, 1.02); paper/sections/03_dynamic.tex:99-106 (R0 calibration); s1_model.tex seq:s1-R0 (peak at m* = sqrt(K_c M_i))

Across the corners the crossings at the assumed dose span 2.41--2.95, 2.54--3.10 and
3.29--3.95 days for 14, 15 and 21 days, and at the capped dose 2.98--3.87, 3.19--4.15 and
4.47--5.80 days. One constant at a time, the 21-day crossing at the assumed dose passes 3.6 days
when $K_c\le50$\uM\ (3.64--3.74 days) or $M_i=150$\uM\ (3.66 days). The fibroblast value of
3.6 days clears the 21-day crossing at six of the eight corners, and 4.3 days clears it at all
eight. The spread over the corners is a range over assumed values, not a probability: nothing
measured puts a distribution on these constants.
% src: audit/brief_2026-10-01/out/t53_crossing_shape.txt:77-79 (D = 6000 ranges; 3.6 d at 6 of 8, 4.3 d at 8 of 8), :89-91 (capped ranges), :10-11,16 (one at a time: 3.735, 3.636, 3.657 d); REPORT.md:1159 (the share of sampled sets is not a probability)

\begin{table}[tbp]
\centering
\small
\caption{\label{stab:s6-corners}Intrinsic critical capacity and crossing half-lives (days) at
the eight corners of the assumed ranges of the carnitine constants, with $R_0$ re-derived at
each corner, and at the nominal constants. Assumed dose: $D=6000$\um\ per fortnight, scaled to
the same amount per day (Eq.~\eqref{seq:s6-cross}). Capped: the dose lowered at each interval and
half-life so that the peak of a patient at 5\% of normal stays at normal capacity
(Eq.~\eqref{seq:s6-capcross}). An interval is safe for every patient, whatever $\theta$, when the
hepatic half-life exceeds the entry.}
\begin{tabular}{@{}rrrcccccccc@{}}
\toprule
& & & & \multicolumn{3}{c}{assumed dose} & & \multicolumn{3}{c}{capped at the ceiling} \\
\cmidrule(lr){5-7}\cmidrule(l){9-11}
$K_c$ (\textmu M) & $M_i$ (\textmu M) & $\rho_0$ & $\thci$ (\um) & 14\,d & 15\,d & 21\,d & & 14\,d & 15\,d & 21\,d \\
\midrule
25 & 150 & 0.35 & 218.2 & 2.90 & 3.04 & 3.89 & & 3.78 & 4.05 & 5.66 \\
25 & 150 & 1 & 232.1 & 2.95 & 3.10 & 3.95 & & 3.87 & 4.15 & 5.80 \\
25 & 1200 & 0.35 & 139.8 & 2.57 & 2.70 & 3.48 & & 3.22 & 3.45 & 4.83 \\
25 & 1200 & 1 & 145.5 & 2.59 & 2.73 & 3.52 & & 3.26 & 3.50 & 4.89 \\
800 & 150 & 0.35 & 127.1 & 2.50 & 2.64 & 3.41 & & 3.12 & 3.34 & 4.68 \\
800 & 150 & 1 & 127.3 & 2.51 & 2.64 & 3.41 & & 3.12 & 3.34 & 4.68 \\
800 & 1200 & 0.35 & 109.9 & 2.41 & 2.54 & 3.29 & & 2.98 & 3.19 & 4.47 \\
800 & 1200 & 1 & 110.0 & 2.42 & 2.55 & 3.30 & & 2.98 & 3.19 & 4.47 \\
\addlinespace
\multicolumn{3}{@{}l}{nominal (100, 300, 0.5)} & 150.8 & 2.62 & 2.75 & 3.55 & & 3.30 & 3.54 & 4.95 \\
\bottomrule
\end{tabular}
\end{table}
% src: audit/brief_2026-10-01/out/t53_crossing_shape.txt:67-75 (corners: theta_crit and D = 6000 crossings), :89-91 (capped nominal 3.30/3.54/4.95 and ranges 2.98-3.87, 3.19-4.15, 4.47-5.80). Capped per-corner entries: eq. seq:s6-capcross (t53:87) evaluated at each corner's theta_crit from model_C.run.theta_crit_intrinsic (R0 re-derived), as t53 does for its ranges; extremes reproduce t53:89-91.

% -----------------------------------------------------------------------------
\subsection{Evidence on the half-life}
\label{si:s6:halflife}

No hepatic half-life of PCC has been measured in humans. Supplementary
Table~\ref{stab:s6-halflife} collects what bears on it. The one human measurement is 3.6 days,
from a [$^{3}$H]biotin pulse--chase in cultured fibroblasts (degradation rate
0.19\,d$^{-1}$)~\cite{chandler1985}. PCC is an enzyme of the mitochondrial matrix, so for the liver
the apt comparator is the median turnover of mouse liver mitochondrial proteins, 4.26
days~\cite{kim2012}, rather than the whole-liver median, 2.4--3.6 days across labelling
studies~\cite{rolfs2021}, which is dominated by faster-turning compartments. Studies of the
delivered enzyme in mouse liver give 5.3--6.7 days for human PCC~\cite{jiang2020dual}.
% src: audit/brief_2026-10-01/pi_note/pi_note.tex:382-390 (3.6 d; matrix enzyme; 4.26 d; whole liver 2.4-3.6 d; 3.8 d intrinsic, about 7 d terminal, 3-5 d; 5.3-6.7 d); REPORT.md:1004 (k_deg = 0.19/d, [3H]biotin pulse-chase); out/t50_hepatic_mito.txt:5-8

The translational pharmacokinetic/pharmacodynamic model of mRNA-3927~\cite{pkpd_translational}
does not supply a single half-life. In it the PCC protein sits in a two-compartment layer. Its
degradation constant, 0.0075\,h$^{-1}$, corresponds to 3.8 days if read as a single
exponential, while the model's stated effective terminal half-life is about 7 days. The local
half-life of that layer is therefore not constant: it passes 3.5 days about five days after a
dose and reaches 9.0 days by the end of a fortnightly interval. What matters for the trough
condition is the delivered capacity at the end of a fortnightly interval, which is reached
before the terminal phase dominates. Matching the periodic trough of the two-compartment layer
with a single exponential gives 4.8 or 3.1 days under two reconstructions of its unpublished
equations (elimination from the central compartment only, or from both), so the translational
model points to about 3--5 days rather than 7.
% src: src/common/params.py:283-299 (k_deg = 0.0075/h; two-compartment; ~7 d; local half-life passes 3.5 d at about five days after a dose and reaches 9.0 d by the end of a fortnightly interval); audit/brief_2026-10-01/out/t48_lambda_literature.txt:5-30 (3.8 d intrinsic; terminal 11.8 / 3.9 d under the two reconstructions; trough-matched 4.8 / 3.1 d; equations unpublished)

From these comparators we take about 4.3 days as the likely hepatic value, plausibly 3--6
days. It rests on comparators, not on a measurement, and its role here is to locate the
question, not to settle it. The sponsor's population model~\cite{gao2025acop} was re-estimated
on human data with an indirect-response structure and may contain a human protein-decay term,
but its parameters are not public; its abstract reports only the clearance (0.266\,L/h) and
volume (13.9\,L) of the mRNA.
% src: audit/brief_2026-10-01/pi_note/pi_note.tex:389-391 (about 4.3 d, plausibly 3-6 d), :403-406 (Gao: re-estimated on human data, indirect response, parameters not public); REPORT.md:1076 (CL = 0.266 L/h, V = 13.9 L; turnover term not public)

\begin{table}[tbp]
\centering
\footnotesize
\caption{\label{stab:s6-halflife}Evidence on the half-life of the delivered PCC protein. None
of these is a measurement in human liver; the likely hepatic value, about 4.3\,d (plausibly
3--6\,d), is an estimate from the comparators.}
\begin{tabular}{@{}p{0.2\linewidth}p{0.3\linewidth}p{0.15\linewidth}p{0.27\linewidth}@{}}
\toprule
source & system and method & half-life (d) & bearing on the hepatic value \\
\midrule
Chandler and Ballard~\cite{chandler1985} & human fibroblast lines; [$^{3}$H]biotin pulse--chase,
degradation rate 0.19\,d$^{-1}$ & 3.6 & the only human measurement of PCC turnover; not liver \\
Kim et al.~\cite{kim2012} & mouse liver mitochondria; metabolic labelling, proteome median &
4.26 & the compartment of PCC, a matrix enzyme \\
Rolfs et al.~\cite{rolfs2021} and other labelling studies & mouse whole liver; proteome median &
2.4--3.6 & includes faster-turning compartments; a less apt comparator \\
Attarwala et al.~\cite{pkpd_translational} & translational PK/PD model of mRNA-3927, fitted in
mice, rats and monkeys; two-compartment protein layer & 3.8 (intrinsic); about 7 (terminal);
3.1--4.8 (trough-matched) & not a single exponential; the trough at fortnightly spacing sees
about 3--5\,d \\
Jiang et al.~\cite{jiang2020dual} & delivered human PCC in mouse liver & 5.3--6.7 & the delivered
enzyme itself, in mouse \\
Gao et al.~\cite{gao2025acop} & population model re-estimated on trial data; indirect-response
structure & not public & may contain a human protein-decay term \\
\bottomrule
\end{tabular}
\end{table}
% src: audit/brief_2026-10-01/pi_note/pi_note.tex:382-406; REPORT.md:1004-1008 (Chandler k_deg 0.19/d; Attarwala 0.0075/h, 3.8 d, ~7 d; Jiang 5.3-6.7 d), :1064 (matrix enzyme; 4.26 d; whole liver 2.4-3.6 d), :1076 (Gao); out/t48_lambda_literature.txt:12-18 (3.1 / 4.8 d)

\paragraph{The specification at 4.3 days.} At the assumed dose all three intervals are then safe
for every patient whatever the carnitine constants, since the highest crossing over the corners
is 3.95 days; at the nominal constants the 21-day margin is about $+165$\um. At the capped dose
the 14- and 15-day intervals remain safe, narrowly at the least favourable constants (crossings
3.87 and 4.15 days; margins of about $+150$ and $+100$\um\ at the nominal constants). The 21-day
interval at the capped dose is not safe for every patient at any corner (crossings 4.47--5.80
days): at the nominal constants the trough is 97\um, 55\um\ short of $\thci$, so a patient would
need about 1.8\% of normal capacity of their own, which the record cannot confirm.
% src: audit/brief_2026-10-01/pi_note/pi_note.tex:390-397 (4.3 d: safe at the booked increment; capped 14/15 narrowly, 3.87 and 4.15 d; 21 d 4.47-5.80 d; about 1.8 %); figure_guide.tex:506-517 (+165; -55; 97; 1.8 %; +150 and +100)

% -----------------------------------------------------------------------------
\subsection{Uncertainty in the half-life}
\label{si:s6:prob}

Because 4.3 days rests on comparators, we also treated the hepatic half-life as log-normal
about 4.3 days, with a standard deviation of the log of 0.35 (with which the plausible range of
3--6 days spans about one standard deviation either side of 4.3 days) or 0.53. The probability
that it exceeds a crossing $t_{1/2}^{\ast}$ is $\Phi\bigl(\ln(4.3/t_{1/2}^{\ast})/s\bigr)$, with $s$
the standard deviation of the log and $\Phi$ the standard normal distribution function
(Supplementary Table~\ref{stab:s6-prob}). At the assumed dose the nominal 21-day crossing is
cleared with probability 64--71\%, and the crossing at the least favourable constants with
56--60\%. At the capped dose the 21-day crossing (4.95 days) lies above 4.3 days, so the
probability falls to 34--39\%, and here the wider distribution gives the higher value.
% src: audit/brief_2026-10-01/out/t53_crossing_shape.txt:80-81 (64/71 %; 56/60 % at 3.95 d), :89-91 (capped: 69/77, 64/71, 39/34 %); out/t50_hepatic_mito.txt:12-14 (complement: 36 % / 29 % below 3.5452 d); pi_note/pi_note.tex:397-401

\begin{table}[tbp]
\centering
\small
\caption{\label{stab:s6-prob}Probability that the hepatic half-life exceeds a crossing, with
the half-life log-normal about 4.3\,d. Least favourable: the corner $K_c=25$\uM, $M_i=150$\uM,
$\rho_0=1$ of the assumed carnitine ranges.}
\begin{tabular}{@{}llccc@{}}
\toprule
interval and dose & constants & crossing (d) & $s=0.53$ & $s=0.35$ \\
\midrule
21 days, assumed dose & nominal & 3.55 & 64\% & 71\% \\
21 days, assumed dose & least favourable & 3.95 & 56\% & 60\% \\
14 days, capped & nominal & 3.30 & 69\% & 77\% \\
15 days, capped & nominal & 3.54 & 64\% & 71\% \\
21 days, capped & nominal & 4.95 & 39\% & 34\% \\
\bottomrule
\end{tabular}
\end{table}
% src: audit/brief_2026-10-01/out/t53_crossing_shape.txt:80-81,89-91

% -----------------------------------------------------------------------------
\subsection{A threshold that allows for illness}
\label{si:s6:illness}

The intrinsic threshold describes a patient at baseline load. A stricter threshold can be
defined in the linear-clearance model, which has no tipping point, as the capacity at which an
illness load holds the marker at a tolerable level (Eq.~\eqref{seq:s1-thcritlin}). With the
illness load at 1.5 times that model's baseline load of 323\um\ and a tolerable level of
40\uM, it is 994\um. The fortnightly trough exceeds it only for half-lives above about 5.0 days
(it is 400\um\ at 3.5 days and 702\um\ at 4.3 days). Against this threshold, Eq.~\eqref{seq:s6-cross}
gives crossings at the assumed dose of about 5.0, 5.2 and 6.3 days for 14, 15 and 21 days, so at
the likely 4.3 days every interval again depends on $\theta$.
% src: audit/brief_2026-10-01/out/t44_interval_decision.txt:40 (theta_crit = 994.3 at fold 1.5); pi_note/pi_note.tex:426-434 (994; above Dtr(14) below about 5.0 d; 400 and 702; crossings about 5.0, 5.2, 6.3 d); s1_model.tex seq:s1-thcritlin (430 at f = 1, 994 at f = 1.5)

We report this threshold as a qualification, not as the specification. It is taken from a
different model from the one that defines the trough condition, its tolerable level is a
convention, and its value moves with the assumed illness amplitude (430\um\ without illness,
994\um\ at 1.5 times the load). The intrinsic threshold, by contrast, follows from the carnitine
model alone and needs no tolerable level.
% src: s1_model.tex:198-215 (430 at f = 1, 994 at f = 1.5; threshold by convention); audit/brief_2026-10-01/out/t44_interval_decision.txt:30-40 (s_pk/m_tol convention; 429.5 and 994.3)

% -----------------------------------------------------------------------------
\subsection{Estimating \texorpdfstring{$\theta$}{theta} instead}
\label{si:s6:theta}

When the half-life lies below a crossing, the alternative is to estimate each patient's $\theta$
and lengthen the interval only for patients whose estimated margin
$\hat\mu(\tau)=\hat\theta+\dtr(\tau;\hat\lambda)-\thci$ is confidently positive. We tested this in
two ways, both at 21 days and with $\lambda$ estimated from the same record, as it would be in
practice.

By simulation, with nonlinear least-squares fits of $(\theta,K_m,b,k_e,\lambda)$ to records of
weekly draws on therapy at a 3.5-day half-life (the routine record alone), the standard deviation
of $\hat\mu(21)$ stays between 109 and 162\um, comparable to $\thci$ itself, from six months to five
years of draws, for patients at 20, 65 and 150\um. It does not shrink with record length, and the
5th percentile of $\hat\mu$ stays below zero at every capacity and length, so a one-sided lower
bound never clears zero. The reason is that $\hat\theta$ is heavy-tailed on the natural scale at
these capacities: a one-sided bound on the margin is set by the tail of the estimate, which does
not tighten as its log-scale spread falls.
% src: audit/brief_2026-10-01/out/t46_timing.txt:12-34 (NLS mu_hat: sd 141-162 at 150, 121-139 at 65, 109-126 at 20; p5 below 0 throughout; does not tighten); REPORT.md:976-978 (one-sided bound is tail-dominated; read natural-scale decisions by simulation)

With both tracer measurements of Supplementary Section~\ref{si:s4} (the intravenous tracer, giving
$k_e$ to 5.4\%, and the isotope-dilution measurement of the load turnover, to 10\% of the load)
added to the 25 monthly draws of the recommended design, Fisher information gives a standard
deviation of the 21-day margin of 92--95\um\ at every capacity from 0.17 to 5\% of normal, and still
74--79\um\ with $k_e$ and the unrecorded load known exactly. We applied a one-sided 95\% rule
(lengthen when $\hat\mu-1.645\,\operatorname{sd}(\hat\mu)>0$) to a population with $\theta$
log-uniform between 0.17 and 5\% of normal (Supplementary Table~\ref{stab:s6-certify}). At a
3.5-day half-life and the nominal constants, where the 21-day interval is in truth safe for 87\% of
this population, the tracers certify 12\% of patients. One patient's estimate certifies the
interval half the time only at about 5.5\% of normal, and 80\% of the time at about 8.2\%, at or
above the top of the range in which the question arises. A check with simulated fits agreed with
these Fisher-information values to within about one to two percentage points. The routine record
without tracers is not included, because without them $K_m$ is unconstrained and the prior on
$k_e$ is wide, so the linear-Gaussian approximation does not hold for it.
% src: audit/brief_2026-10-01/out/t54_dose_timing.txt:75-107 (margin sd 92-95 and 74-79; certification rule; population; 87 %, 12 %, 5.5 % / 8.2 %; record alone not shown); t54_dose_timing.py:24-30 (log-uniform 0.17-5 %; simulated-fit check within about 1-2 points), :213 (tracer 0.054, load turnover 0.10); si/s4_routes.tex:163-165 (10 %, 5.4 %)

\begin{table}[tbp]
\centering
\footnotesize
\caption{\label{stab:s6-certify}Certifying a 21-day interval patient by patient from an estimate
of $\theta$, with $\lambda$ estimated from the same record (25 monthly draws plus the stated
measurements; Fisher information, one-sided 95\% rule). Population: $\theta$ log-uniform over
0.17--5\% of normal. $\theta_{\mathrm{need}}=\thci-\dtr(21)$; truly safe: share of the population
with $\theta>\theta_{\mathrm{need}}$. Certified: share of all patients; of the safe: share of the
truly safe; false-safe: share of the truly unsafe who are certified. Last column: capacity at
which one patient's estimate certifies the interval half the time and 80\% of the time. Tracers:
$k_e$ to 5.4\% and the load turnover to 10\%. Exact: $k_e$ and the unrecorded load known. The
third setting scales $D=2672$ to the same amount per day, so its 21-day peak exceeds the ceiling.}
\begin{tabular}{@{}p{0.22\linewidth}llcccc@{}}
\toprule
setting ($\theta_{\mathrm{need}}$; truly safe) & measurements & certified & of the safe &
false-safe & capacity, 50\% / 80\% \\
\midrule
3.5\,d, $D=6000$, nominal & tracers & 12\% & 13\% & 4.8\% & 5.5\% / 8.2\% of normal \\
(7.9\um; 87\%) & exact & 14\% & 16\% & 4.7\% & 4.6\% / 6.8\% \\
\addlinespace
3.5\,d, $D=6000$, least favourable & tracers & 3\% & 9\% & 1.3\% & 8.2\% / 10.9\% \\
(89.2\um; 15\%) & exact & 2\% & 10\% & 1.0\% & 7.3\% / 9.5\% \\
\addlinespace
4.3\,d, $D=2672$ scaled, nominal & tracers & 11\% & 13\% & 4.7\% & 5.9\% / 8.7\% \\
(10.3\um; 79\%) & exact & 13\% & 15\% & 4.7\% & 5.0\% / 7.5\% \\
\bottomrule
\end{tabular}
\end{table}
% src: audit/brief_2026-10-01/out/t54_dose_timing.txt:94-107

% -----------------------------------------------------------------------------
\subsection{Reading the half-life from the dated record}
\label{si:s6:record}

The half-life, unlike $\theta$, is identified by the routine record on therapy, because the
marker's response to each dated infusion depends on $\lambda$ directly. We simulated two years of
the record of a patient at 5\% of normal, with the irregularities of real dosing: infusion dates
jittered with a standard deviation of 2 days, 8\% of infusions missed and one 28-day supply gap.
With $\ln\lambda$ as a fifth unknown alongside $\theta$, $K_m$, $b$ and $k_e$, Fisher information
gives $\sdl{\lambda}=0.055$ with monthly draws (25 over two years), 0.046 fortnightly (52), 0.037
weekly (104) and 0.023 on alternate days (365). The irregularities also help $\theta$
($\sdl{\hat\theta}$ falls from 3.05 for a regular train to 2.28 with all three, monthly), but
freeing $\lambda$ costs precision on $\theta$, from 2.28 to 3.45 with monthly draws (1.69 to 3.01
fortnightly, 1.27 to 2.45 weekly, 0.75 to 1.61 on alternate days), because the two share part of
the response to each dose. Since $\theta$ is not identified in severe PA either way, little that
the decision could use is given up.
% src: audit/brief_2026-10-01/out/t45_on_therapy.txt:14-18 (3.05 -> 2.90 -> 2.45 -> 2.28; jitter sd 2 d, 8 % missed, 28 d gap), :45-50 (sd ln lambda 0.055/0.046/0.037/0.023; theta 2.28->3.45, 1.69->3.01, 1.27->2.45, 0.75->1.61); REPORT.md:924-934

Over the length of the record, $\sdl{\lambda}$ is 0.094, 0.067 and 0.053 after six months of
monthly, fortnightly and weekly draws, and 0.055--0.058 after two years of monthly draws,
depending on the simulated record. A precision of 25\%, about the distance from 3.5 days to the
nominal 14-day crossing, is reached by six months at every cadence; 5\% takes 48, 24 and 12
months at monthly, fortnightly and weekly draws. Six months of monthly draws already separate a
half-life of 3.6 days from one of 7 days by 7 standard deviations.
% src: audit/brief_2026-10-01/out/t46_timing.txt:42-52 (0.094/0.067/0.053 at 6 months; 0.058 at 24 months monthly; 25 % by 6 mo; 5 % by 48/24/12 mo); out/t45_on_therapy.txt:47 (0.055); pi_note/pi_note.tex:408-410 (0.094; 0.055-0.058; 7 sigma); out/t48_lambda_literature.txt:94 (d' = 7.0 at 6 mo monthly); out/t44_interval_decision.txt:150-153 (25 % margin between 3.5 d and 2.62 d)

Precision on $\lambda$ is not the same as placing $\lambda$ on one side of a crossing. A half-life of
3.6 days lies $\ln(3.6/3.5452)=0.015$ from the nominal 21-day crossing, 43 times closer than 7 days
is to 3.6 days ($\ln(7/3.6)=0.665$). Supplementary Table~\ref{stab:s6-record} gives the expected
separation $z=\lvert\ln(t_{1/2}/t_{1/2}^{\ast})\rvert/\sdl{\lambda}$, with the standard deviation
evaluated at the true half-life, and the chance $\Phi(z-2)$ that one noisy record places the value
on the correct side at two standard deviations. One patient's record never places a 3.6-day truth
on a side of the nominal crossing: at best it sits 0.58 standard deviations away after five years of
weekly draws, inside a two-standard-deviation band around the crossing that is never narrower than
3.36--3.74 days. A 4.3-day truth is placed above the nominal crossing with 83\% probability after a
year of monthly draws, but above the highest corner crossing at the assumed dose (3.95 days) with
only 24\% (88\% after five years of weekly draws).
% src: audit/brief_2026-10-01/out/t52_knife_resolution.txt:5-7 (0.665; 0.0153; 43x), :12-30 (z, P, dead zones), :32-34 (best 0.58 sigma), :47,63 (vs 3.95 d: 24 % at 12 mo monthly; 88 % at 60 mo weekly); pi_note/pi_note.tex:411-419

Pooling patients helps only if they share one half-life. If they do, a 3.6-day truth could be placed
at two standard deviations with 77 patients at 50\% power, or 155 at 80\%, from a year of monthly
draws each; with 35 or 71 from a year of weekly draws; and with 23 or 46 from two years of weekly
draws. These 23--155 patients exceed the sixteen the trial enrolled~\cite{trial_phase12}. Between-patient
spread in hepatic PCC turnover is unmeasured; any spread raises these numbers or turns the pooled
value into a population mean that no single patient has.
% src: audit/brief_2026-10-01/out/t52_knife_resolution.txt:65-68 (77/155, 35/71, 23/46); REPORT.md:1435 (pooling assumes identical lambda; spread raises counts or answers a population mean); manuscript/main.tex:431-432 (sixteen enrolled)

These record-based numbers are Fisher-information values, not simulated fits. The Gaussian
approximation for $\ln\lambda$ was checked at 3.5 days, and the chances $\Phi(z-2)$ are one-sided
decisions that we have not simulated. Two further routes could supply the hepatic value: the protein
decay term of the sponsor's population model~\cite{gao2025acop}, if released, and a direct measurement
of PCC turnover in human hepatocytes by pulse--chase or dynamic SILAC, with the caveat that cultured
cells can turn proteins over differently from the liver in vivo. For the delivered increment, the
paired breath tests of the main text's data specification establish the same-day repeatability on
which a measurement of $D$ before and after an infusion depends. Dated delays in the record, the
subject of Supplementary Section~\ref{si:s7}, are part of what identifies $\lambda$.
% src: REPORT.md:1435 (FIM, not simulated; Gaussian regime confirmed at 3.5 d only; Phi(z-2) not simulated); pi_note/pi_note.tex:404-406,420-425 (Gao; hepatocyte pulse-chase or dynamic SILAC, caveat; Tier 2 breath tests for D), :543-548 (a late infusion is also a measurement)

\begin{table}[tbp]
\centering
\footnotesize
\caption{\label{stab:s6-record}Can one patient's record place the half-life on a side of the
21-day crossing? Fisher information for $(\ln\theta,\ln K_m,b,\ln k_e,\ln\lambda)$ on the
on-therapy record of a patient at 5\% of normal, with dosing irregularities (jitter, missed
infusions, a supply gap), re-evaluated at each true half-life. $z$: expected separation from the
crossing in standard deviations. $P$: chance $\Phi(z-2)$ that one record places a 4.3-day truth
above the crossing at two standard deviations, against the nominal crossing (3.55\,d) and the
highest corner crossing at the assumed dose (3.95\,d). Last column: true half-lives that one record
cannot place on either side of the nominal crossing at two standard deviations.}
\begin{tabular}{@{}lcccccccc@{}}
\toprule
& & \multicolumn{2}{c}{truth 3.6\,d} & & \multicolumn{3}{c}{truth 4.3\,d} & \\
\cmidrule(lr){3-4}\cmidrule(lr){6-8}
cadence & months & $\sdl{\lambda}$ & $z$ & & $z$ (3.55) & $P$ (3.55) & $P$ (3.95) & band (d) \\
\midrule
monthly & 6 & 0.095 & 0.16 & & 1.93 & 47\% & 12\% & 2.93--4.28 \\
monthly & 12 & 0.067 & 0.23 & & 2.97 & 83\% & 24\% & 3.10--4.05 \\
monthly & 24 & 0.057 & 0.27 & & 3.47 & 93\% & 32\% & 3.16--3.98 \\
monthly & 60 & 0.045 & 0.34 & & 4.36 & 99\% & 47\% & 3.24--3.88 \\
fortnightly & 6 & 0.067 & 0.23 & & 2.92 & 82\% & 24\% & 3.10--4.05 \\
fortnightly & 24 & 0.047 & 0.33 & & 4.18 & 99\% & 44\% & 3.23--3.89 \\
weekly & 6 & 0.053 & 0.29 & & 3.71 & 96\% & 36\% & 3.19--3.94 \\
weekly & 12 & 0.045 & 0.34 & & 4.43 & 99\% & 48\% & 3.24--3.88 \\
weekly & 24 & 0.037 & 0.42 & & 5.27 & 100\% & 63\% & 3.30--3.81 \\
weekly & 60 & 0.027 & 0.58 & & 7.22 & 100\% & 88\% & 3.36--3.74 \\
\bottomrule
\end{tabular}
\end{table}
% src: audit/brief_2026-10-01/out/t52_knife_resolution.txt:12-30 (sd at 3.6 d, z at 3.6 d, z and P at 4.3 d, dead zones), :46-63 (P vs 3.95 d)
